% Auto-generated by paper/audit/behavior_stratified.py -- do not edit by hand.
\begin{table*}[t]
\centering
\small
\begin{tabular}{l rr rrrr l}
\toprule
Generator behavior (step 5) & verdicts & answers & TP & FP & FN & TN & Headline rate [95\% CI] \\
\midrule
context-follow (adopts false value) & 1212 & 405 & 1025 & 0 & 187 & 0 & recall 84.6 [81.5, 87.6] \\
both-claims (true + false value) & 135 & 45 & 88 & 0 & 44 & 3 & recall 66.7 [54.5, 78.8] \\
memory-override (uses true value) & 42 & 14 & 0 & 0 & 0 & 42 & FPR 0.0 [0.0, 7.1] \\
conflict-aware (notices source is off) & 54 & 18 & 3 & 46 & 0 & 5 & FPR 90.2 [76.5, 100.0] \\
refusal / insufficient & 400 & 134 & 23 & 130 & 1 & 246 & FPR 34.6 [27.4, 42.2] \\
unrelated / failed & 606 & 202 & 0 & 25 & 0 & 581 & FPR 4.1 [2.1, 6.3] \\
\midrule
unlabeled (no llm\_eval) & 9 & 3 & 0 & 8 & 0 & 1 & FPR 88.9 [66.7, 100.0] \\
\bottomrule
\end{tabular}
\caption{Judge outcomes (step 6) stratified by generator behavior (step 5), validated records, pooled over all nine judge cells. TP/FP/FN/TN are true/false positives/negatives against the induced-error gold. The headline rate is recall for committed answers and false-positive rate (FPR) otherwise, with a 95\% cluster bootstrap CI resampled by answer (10,000 replicates). Recall is concentrated in context-follow answers; false positives are concentrated in conflict-aware and refusal answers, while memory-override controls are essentially never flagged. Behavior labels are the LLM label-evaluation output; the three outputs without a label are shown separately as unlabeled.}
\label{tab:behavior-stratified}
\end{table*}
